\documentclass[a4paper,12pt]{article}
\setcounter{secnumdepth}{5}
\setcounter{tocdepth}{3}
\input{/usr/share/LaTeX-ToolKit/template.tex}
\begin{document}
\title{Special Theory of Relativity}
\author{沈威宇}
\date{\temtoday}
\titletocdoc
\sct{Special Theory of Relativity}
\ssc{Postulates}
The special theory of relativity, or simply special relativity, is a scientific theory of the relationship between space and time. In Albert Einstein's 1905 paper, "On the Electrodynamics of Moving Bodies", the theory is presented as being based on just two postulates:
\bem
\item Principle of Relativity: The laws of physics are invariant (identical) in all inertial frames of reference (that is, frames of reference with no acceleration, that is, frames in which Newton's laws are valid).
\item Principle of Light Constancy or Principle of Light Speed Invariance: The speed of light in vacuum is the same for all observers, regardless of the motion of light source or observer.
\een
\ssc{Terminology}
\bit
\item Speed of light $c$: the maximum speed of information, independent of the speed of the source and receiver.
\item Clock: a device to measure differences in time.
\item Event: something that happens at a definite place and time.
\item Spacetime: geometrical space and time considered together.
\item Spacetime interval between two events: a measure of separation between events that is defined as square their separation in time minus square their separation in geometrical space.
\item Reference frame or coordinate system: a way to locate events in spacetime. Events have coordinates $x,y,z$ for space and $t$ for time.
\item Inertial reference frame or free-float reference frame: a reference frame where objects follow Newton's law.
\item Coordinate transformation: changing how an event is described from one reference frame to another.
\item Invariance: when physical laws or quantities do not change in different inertial frames. The speed of light is invariant in special relativity.
\eit
\ssc{Lorentz transformation}
\sssc{Lorentz transformation and its inverse}
Define an event to have spacetime coordinates $(t,x,y,z)$ in reference frame $S$ and $(t',x',y',z')$ in reference frame $S'$ moving at a velocity $v$ in positive $x$ direction. Then the Lorentz transformation specifies that these coordinates are related in the following way:
\[t'=\gamma\qty(t-\frac{v}{c^2}x)\]
\[x'=\gamma(x-vt)\]
\[y'=y\]
\[z'=z\]
where
\[\gamma=\frac{1}{\sqrt{1-\beta^2}}=\frac{1}{\sqrt{1-\qty(\frac{v}{c})^2}}\]
is the Lorentz factor, $c$ is the speed of light, and $\beta=\frac{v}{c}$ is the speed parameter.

The inverse Lorentz transformation is obtained by solving the four Lorentz transformation equations as
\[t=\gamma\qty(t'+\frac{v}{c^2}x)\]
\[x=\gamma(x'+vt')\]
\[y=y'\]
\[z=z'\]
A quantity that is invariant under Lorentz transformations is known as a Lorentz scalar.

Suppose one event has coordinates $(t_1,x_1,y_1,z_1)$ in $S$ and $(t_1',x_1',y_1,z_1)$ in $S'$ and another event has coordinates $(t_2,x_2,y_2,z_2)$ in $S$ and $(t_2',x_2',y_2,z_2)$ in $S'$. Let
\[\Delta t=t_2-t_1\]
\[\Delta x=x_2-x_1\]
\[\Delta t'=t_2'-t_1'\]
\[\Delta x'=x_2'-x_1'\]
We get
\[\Delta t'=\gamma\qty(\Delta t-\frac{v}{c^2}\Delta x)\]
\[\Delta x'=\gamma(\Delta x-v\Delta t)\]
\[\Delta t=\gamma\qty(\Delta t'+\frac{v}{c^2}\Delta x')\]
\[\Delta x=\gamma(\Delta x'+v\Delta t')\]
If we take differentials instead of differences, we get
\[\dd{t'}=\gamma\qty(\dd{t}-\frac{v}{c^2}\dd{x})\]
\[\dd{x'}=\gamma(\dd{x}-v\dd{t})\]
\[\dd{t}=\gamma\qty(\dd{t'}+\frac{x}{c^2}\dd{x'}\]
\[\dd{x}=\gamma(\dd{x'}+v\dd{t'})\]
\sssc{Derivation of Lorentz transformation}
By the two postulates, we have
\[(ct)^2-x^2-y^2-z^2=(ct')^2-x'^2-y'^2-z'^2.\]
Since the motion is along the $x$-axes only, we have
\[(ct)^2-x^2=(ct')^2-x'^2.\]
By Newton's law, the solution must take the linear form
\[x'=Ax+Bct\]
\[ct'=Cx+Dct\]
Substitution gives
\[c^2t^2-x^2=C^2x^2+2CDcxt+D^2c^2t^2-A^2x^2-2ABcxt-B^2c^2t^2\]
\[-1=C^2-A^2\]
\[A^2-C^2=1\]
\[0=2CDcxt-2ABcxt\]
\[AB=CD\]
\[c^2=D^2c^2-B^2c^2\]
\[D^2-B^2=1\]
\[D=\frac{AB}{C}\]
\[A^2-C^2=\qty(\frac{C}{B})^2=1\]
Take
\[B=C\]
\[A=D\]
Since
\[\cosh^2\psi-\sinh^2\psi=1\]
Let
\[A=D=\cosh\psi\]
\[B=C=\sinh\psi\]
We have
\[x'=\cosh\psi x+\sinh\psi ct\]
\[ct'=\sinh\psi x+\cosh\psi ct\]
Substituting $x'=0$ in $S'$ is measured as $x=vt$ in $S$
\[0=\cosh\psi vt+\sinh\psi ct\]
\[ct'=\cosh\psi ct+\sinh\psi vt\]
Using identities
\[\sinh\psi=\frac{\tanh\psi}{\sqrt{1-\tanh^2\psi}}\]
\[\cosh\psi=\frac{1}{\sqrt{1-\tanh^2\psi}}\]
We have
\[0=vt+\tanh\psi ct\]
\[\tanh\psi=-\frac{v}{c}=-\beta\]
\[\sinh\psi=\frac{-\beta}{\sqrt{1-\beta^2}}=-\gamma\beta\]
\[\cosh\psi=\frac{1}{\sqrt{1-\beta^2}}=\gamma\]
Substitution gives
\[x'=\gamma\qty(x-vt)\]
\[ct'=\gamma(ct-\beta x)\]
\[t'=\gamma\qty(t-\frac{v}{c^2}x)\]
\ssc{Consequence of Lorentz transformation}
\sssc{Invariant interval}
In Galilean relativity, the spatial separation or separation in space $\sqrt{(\Delta x)^2+(\Delta y)^2+(\Delta z)^2}$, and the temporal separation or separation in time, $\Delta t$, between two events are independent invariants. In special relativity, however, they are not, and the invariant interval, denoted as $(\Delta s)^2$ and defined as
\[(\Delta s)^2=c^2(\Delta t)^2-\qty((\Delta x)^2+(\Delta y)^2+(\Delta z)^2)\]
is invariant.
\bpr
Consider the Lorentz transformation setup.
\[\Delta s=c^2(\Delta t)^2-(\Delta x)^2\]
\[\ba
\Delta s'^2=&c^2(\Delta t')^2-(\Delta x')^2\\
=&\gamma^2\qty(\qty(c\Delta t-\beta\Delta x)^2-(\Delta x-v\Delta t)^2)\\
=&\gamma^2\qty(c^2(\Delta t)^2-2v\Delta x\Delta t+\beta^2(\Delta x)^2-(\Delta x)^2+2v\Delta x\Delta t-v^2(\Delta t)^2)\\
=&\gamma^2\qty(\qty(\frac{c}{\gamma})^2(\Delta t)^2-\frac{1}{\gamma^2}(\Delta x)^2)\\
=&c^2(\Delta t)^2-(\Delta x)^2\\
=&(\Delta s)^2
\ea\]
\epr
There are three cases of $(\Delta s)^2$:
\bit
\item $(\Delta s)^2>0$: This implies that $\abs{\frac{\Delta x}{\Delta t}}<c$, called the two events are timelike separated. In this case, it is possible to find a inertial reference frame in which the two events happen at the same place. In this frame, the separation in time is $\frac{\sqrt{(\Delta s)^2}}{c}$, called proper time (interval).
\item $(\Delta s)^2<0$: This implies that $\abs{\frac{\Delta x}{\Delta t}}>c$, called the two events are spacelike separated. In this case, it is possible to find a inertial reference frame in which the two events happen at the same time. In this frame, the separation in space is $\sqrt{-(\Delta s)^2}$, called proper distance or proper length.
\item $(\Delta s)^2=0$: This implies that $\abs{\frac{\Delta x}{\Delta t}}=c$, called the two events are lightlike separated.
\eit
\sssc{Relativity of simultaneity}




\sssc{Time dilation and Langevin's light-clock}
Two events that happen at the same spatial point with time separation $\Delta t$ in $S$ has time separation $\Delta t'$ in $S'$:
\[\Delta t'=\frac{\Delta t}{\sqrt{1-\qty(\frac{v}{c})^2}}=\frac{\Delta t}{1+\beta^2}=\gamma\Delta t.\]
Langevin's light-clock: 

A light-clock is imagined to be a box of perfectly reflecting walls wherein a light signal reflects back and forth from opposite faces.

Consider the scenario that observer A on a train holds a light-clock of $L$ as well as an electronic timer with which she measures how long it takes a pulse to make a round trip up and down along the light-clock. From her point of view the emission and receipt of the pulse occur at the same place, and she measures the interval using a single clock located at the precise position of these two events. For the interval between these two events, observer A finds
\[\Delta t_0=\frac{2L}{c},\]
which is a proper time interval.

Another observer B, who is parked by the tracks as the train goes by at a speed of $v$⁠. Instead of making straight up-and-down motions, observer B sees the pulses moving along a zig-zag line. However, because of the postulate of the constancy of the speed of light, the speed of the pulses along these diagonal lines is the same $c$. For the interval between these two events, observer B finds
\[\Delta t=\frac{2L}{\sqrt{c^2-v^2}}=\frac{\Delta t_0}{\sqrt{1-\qty(\frac{v}{c})^2}},\]
which is not a proper time interval.



\ssc{Length contraction}

